\section{Related Work}
\paragraph{Selection bias and adaptive reuse of evaluation data.}
That the best of many models looks better on the data used to pick it than on new data is the classic selection-bias problem in model assessment \citep{cawley2010overfitting}; bootstrap corrections of such optimism \citep{efron1997632} and nested resampling are standard in the statistical-learning literature.
In machine learning competitions and benchmarks, adaptive reuse of one test set has been studied through the lens of adaptive data analysis: the reusable holdout (Thresholdout) \citep{dwork2015generalization,dwork2015preserving} and the Ladder mechanism \citep{blum2015ladder} limit what a sequence of queries can reveal about the data, and large-scale replication studies found that benchmark accuracies hold up better than feared under test-set reuse but fall under distribution shift \citep{recht2019imagenet}.
Our setting differs in that the ``analyst'' is an automatic tuner whose queries (candidate prompts or programs) are numerous, adaptive, and issued without human restraint, and the reported object is the score of the single winning configuration.

\paragraph{Inference after selection.}
Econometrics and statistics offer estimators and intervals for the \emph{winner} of a selection: conditional and hybrid inference \citep{andrews2024inference}, polyhedral conditioning \citep{lee2016exact}, valid post-selection inference \citep{berk2013valid}, and Gaussian multiplier-bootstrap bounds for maxima \citep{chernozhukov2013gaussian}.
We use these as ingredients and ask how they behave for binary per-item outcomes shared across configurations and for search algorithms (successive halving, surrogate-guided search) whose selection event is not a simple polyhedron.

\paragraph{Prompt and program optimisation, and evaluation noise.}
Automatic search over instructions and demonstrations is now routine \citep{zhou2022large,khattab2023dspy,opsahl2024optimizing}, including budget-aware formulations as best-arm identification \citep{shi2024efficient}.
Prompt-sensitivity studies show that formatting and wording alone move benchmark scores by many points \citep{sclar2023quantifying,mizrahi2024state}, and efficient multi-prompt evaluation estimates the whole performance distribution over templates \citep{polo2024efficient}; those works treat the prompt as given, whereas we study the bias introduced when the prompt is \emph{chosen} on the benchmark.
\citet{perez2021true} pointed out that prior ``few-shot'' work tuned prompts and hyper-parameters on many held-out examples and studied model selection (cross-validation, minimum description length) when no such examples are available; \citet{miller2024adding} advocates error bars for evaluations, but treats the evaluated system as fixed.
